problem:  
Let \((X,d,\mathfrak{m},p)\) be a non‑collapsed \(RCD(K,N)\) space arising as a measured Gromov–Hausdorff limit of smooth \(n\)-dimensional Riemannian manifolds \((M_i,g_i,p_i)\) with  
\(\mathrm{Ric}_{M_i}\ge K\) and \(\operatorname{vol}(B_1(p_i))\ge v_0>0\).  
For \(x\in X\), let \(\Sigma_x\) denote the collection of all tangent cones (in the pointed measured Gromov–Hausdorff sense) at \(x\), and let  
\[
S_{\mathrm{mult}} := \{\,x\in X : |\Sigma_x|>1 \,\},
\]
the set of points where tangent cones are not unique.  
For each \(x\), define the *Euclidean factor dimension*
\[
k(x):=\max\{\,k\in\mathbb{N} : \exists\,T\in\Sigma_x\text{ admitting an isometric }\mathbb{R}^k\text{ factor}\,\}.
\]

---

**Open problem:**  
1. (*Stratified multiplicity*) For \(0\le k\le N\), set  
\[
S^{k}_{\mathrm{mult}} := \{\,x\in S_{\mathrm{mult}} : k(x)=k \,\}.
\]
Does one always have a quantitative stratification bound
\[
\dim_{\mathrm{Haus}}(S^{k}_{\mathrm{mult}}) \le k
\tag{*}
\]
for all \(k\)?  Equivalently: is the set of points where tangent cones vary but share the same maximal Euclidean factor of dimension \(k\) contained in the classical \(k\)‑stratum of Cheeger–Naber type?

2. (*Non‑uniqueness codimension*)  In particular, does \(\dim_{\mathrm{Haus}}(S_{\mathrm{mult}})\le N-2\) hold for all non‑collapsed \(RCD(K,N)\) spaces?  That is, can tangent cone multiplicity occur on a codimension‑1 subset, or is it forced to be rarer, tied intrinsically to the two‑dimensional or higher singular strata?

3. (*Factor‑dimension stability*)  Assuming \((*)\) holds, does it follow that for almost every \(x\in S_{\mathrm{mult}}\), all tangent cones at \(x\) share the same Euclidean factor \(\mathbb{R}^{k(x)}\), even if their residual cross‑sections differ?

---

Why is it a "good" problem:  
This problem is a refined and conceptually motivated investigation into one of the remaining frontiers in Ricci limit and \(RCD\) geometry: the locus of points with multiple tangent cones. Current theories (Cheeger–Colding, Cheeger–Naber, Mondino–Naber) describe rectifiability and dimensional bounds for singular strata defined via *symmetry loss*, but do not relate those strata directly to *tangent cone multiplicity*. The proposed question bridges these notions, asking whether non‑uniqueness itself is stratified and therefore constrained by the same dimension estimates as loss of symmetries.  

This formulation not only captures a deep geometric phenomenon—how infinitesimal data can bifurcate or remain coherent—but also provides a natural extension of quantitative stratification to the non‑smooth setting. The conjectured codimension‑two limitation is no longer arbitrary: it expresses the hypothesis that multiple tangent cones arise only once at least two independent symmetries fail, connecting multiplicity to analytic degeneracy of harmonic splitting maps and thereby giving precise geometric intuition.  

A resolution would unify the analytic stratification theory, measure‑theoretic rectifiability, and the infinitesimal geometry of Ricci limits, potentially revealing new rigidity or instability mechanisms in synthetic curvature spaces. Simple to state yet demanding to resolve, it epitomizes the “narrow entrance, profound depth’’ ideal of a good, foundational research problem.